\documentclass[10pt,letterpaper]{article}
\usepackage{technical-testing}
\setDocTitle{Python Rational Time And Exposure Sheet}
\setDocSubtitle{Implementation laboratory with reproducibility contract}
\setDocVersion{1.0}
\setDocStatus{Educational / Draft}
\setDocOwner{Game Development}
\setDocDate{2026-10-30}
\setDocSummary{Implementation laboratory with reproducibility contract}
\begin{document}
\makedoctitle
\makedocsummary

\section{Prerequisites and source crosswalk}
Python 3.11; fractions; CSV. Book chapters: Chapters 9--17, 31--33. Engineering procedure and code are supplementary.
\section{Input and baseline}
FrameRate, ExposureSheet, on\_twos in code/python/animation\_time.py. Record shot ID, coordinate units, target fps, time convention, compiler/interpreter and fixture revision.
\section{Experiment procedure}
\begin{enumerate}
\item Generate 60 display-frame rows at 24 fps; add three nonoverlapping cues; inspect end boundary at frame 60.
\item Run a baseline and two controlled variations; change only one parameter per variation.
\item Capture source data (CSV or trace), a sampled frame strip, and test command output.
\item Identify at least one scenario where the intended visual effect conflicts with physical realism or performance limits.
\end{enumerate}
\section{Acceptance criteria}
Run test\_animation.py; verify 60 rows, 30 drawings, no gaps; introduce an overlapping exposure and confirm rejection. No result is considered verified without stored output and test command.
\section{Required evidence bundle}
\begin{itemize}
\item Written hypothesis and expected perceptual result.
\item Environment/configuration manifest and exact execution commands.
\item Raw outputs, plots or contact sheets, comparison table and observed failure modes.
\item Engineering decisions, reviewers, and one prioritized revision.
\end{itemize}
\end{document}
